\section{Guidance}\label{sec:guidance}

The rules below follow from the results above. Each names the evidence it rests on.

\begin{enumerate}[label=\textbf{Rule \arabic*.},leftmargin=*,align=left]
  \item \emph{Count the positives in the data the selector sees.} The minority count, not the prevalence, sets what can be found (\cref{thm:limit}). With cross-validation or a held-out split, count the positives in the training part only, and run the selection inside each fold, so that the model's performance is estimated on data the selection never saw.
  \item \emph{Translate $\nmin$ into the smallest effect that can be found.} A feature's chance of beating the shadows is governed by $s = \nmin\delta^2$, or by $\sauc = 4\pi\nmin(\AUC - \frac12)^2$ for a feature of any distribution (\cref{sec:rank}). \Cref{tab:sbands} gives the false negative rate of the one-shot rule with five shadows in E2; with more shadows it is higher (\cref{cor:shadows}). A feature with single-feature ROC area 0.6 has $\sauc \approx 0.13\,\nmin$, so it is found reliably only with several hundred positives.
  \item \emph{With XGBoost before 3.0, set the intercept to the prevalence} (\code{base\_score = y.mean()}). The default of versions 1.x and 2.x adds a toll of order $\lambda/\phat$ to every split (\cref{prop:toll}) and raised the false negative rate by \numGapMeanSmall{} on average at $\phat \le 0.01$ in E2.
  \item \emph{Use a fixed, moderate number of shadows.} The level is $1/(K+1)$ whatever the imbalance (\cref{prop:level}), and more shadows cost power (\cref{cor:shadows}). Ten shadows give a level of 0.09 per one-shot comparison; one shadow per candidate, as Boruta does, makes wide tables harder than narrow ones. For discrete columns use per-feature shadows, which keep the level for any column type.
  \item \emph{Screen with a rank test.} At the same level, the Mann--Whitney test lost far fewer relevant features than any tree importance when the signal shifts the class distributions (\cref{tab:methods}), and its power is predicted by $\sauc$ (\cref{prop:mw}). It misses features that act only through their spread or through interactions (\cref{sec:robustness}), so use it to decide what is clearly in, not to drop what a model finds useful.
  \item \emph{With repeated rounds, keep the tentative features, and know the price.} Below about 100 positives, keeping only accepted features drops most moderately predictive ones (0.53 at $\nmin = 50$, $\delta = 0.5$). Keeping tentative features loses few (0.07) but keeps irrelevant candidates at a rate that depends on the number of rounds $R$, the number of shadows $K$ and the number of candidates $d$ (\cref{prop:rounds}); compute it for your $(R, K, d)$ before trusting the selection.
  \item \emph{Do not expect class weights, downsampling or early stopping to help.} None changes $\nmin$, and none helped in E4.
  \item \emph{Below about 25 positives, do not prune on importance.} At $\nmin = 10$ every method we tested either lost most relevant features or kept many irrelevant ones, and at $\nmin = 25$ the best of them, the rank test with Benjamini--Hochberg at $q = 0.2$, still lost 0.24 of the relevant features while keeping 0.15 of the irrelevant ones (\cref{tab:methods}). Keep features on domain grounds, inspect a penalised model's regularisation path, or collect more positives: doubling $\nmin$ helps as much as a feature with a $\sqrt2$ larger effect.
\end{enumerate}

\begin{table}[t]
\centering\small
\caption{False negative rate of the one-shot max-shadow rule ($K = 5$, 300-tree XGBoost, prevalence intercept) by band of $s = \nmin\delta^2$, over the E2 cells in each band.}\label{tab:sbands}
\begin{tabular}{@{}lcccl@{}}
\toprule
$s$ & range over cells & median & cells & example\\
\midrule
\csname @@input\endcsname generated/table_sbands
\bottomrule
\end{tabular}
\end{table}

\paragraph{The package.} \pkg{} implements these rules as scikit-learn compatible estimators. \code{assess\_imbalance} reports $\nmin$ and the smallest detectable effect; \code{UnivariateScreen} runs the rank-test screen with Benjamini--Hochberg control; and \code{ShadowFeatureSelector} runs Boruta-style rounds with any classifier, with ten shadows drawn afresh each round, tentative features kept, a Bonferroni or Benjamini--Hochberg correction, and the prevalence intercept set automatically for XGBoost before 3.0. It warns when the training data hold fewer than 100 positives, refuses a number of rounds too small for any feature to reach significance, and reports each feature's estimated $s$ and the expected keep rate of an irrelevant candidate for the chosen $(R, K, d)$.
